\documentclass[a4paper]{amsart} 
\usepackage{amssymb,times, amscd,amsmath,amsthm, xypic}
\usepackage{tikz}
\usepackage{tikz-cd}
\usepackage[all]{xy}
\usepackage{geometry}
\usepackage{mathrsfs}

\begin{document} 

\numberwithin{equation}{section}
\newtheorem{thm}[equation]{Theorem}
\newtheorem{pro}[equation]{Proposition}
%\newtheorem{prob}[equation]{Problem}
%\newtheorem{qu}[equation]{Question}
\newtheorem{qu}{Question}
\newtheorem{cor}[equation]{Corollary}
\newtheorem{con}[equation]{Conjecture}
\newtheorem{lem}[equation]{Lemma}
\theoremstyle{definition}
\newtheorem{ex}[equation]{Example}
\newtheorem{defn}[equation]{Definition}
\newtheorem{observation}[equation]{Observation}
\newtheorem{rem}[equation]{Remark}

\renewcommand{\rmdefault}{ptm}

\def \calB{\mathcal B}
\def\frak{\mathfrak}
\def\alp{\alpha}
\def\be{\beta}
\def\jeden{1\hskip-3.5pt1}
\def\om{\omega}
\def\bigstar{\mathbf{\star}}
\def\ep{\epsilon}
\def\vep{\varepsilon}
\def\Om{\Omega}
\def\la{\lambda}
\def\La{\Lambda}
\def\si{\sigma}
\def\Si{\Sigma}
\def\Cal{\mathcal}
\def\m {\mathcal}
\def\ga{\gamma}
\def\Ga{\Gamma}
\def\de{\delta}
\def\De{\Delta}
\def\bF{\mathbb{F}}
\def\bH{\mathbb H}
\def\bPH{\mathbb {PH}}
\def \bB{\mathbb B}
\def \bA{\mathbb A}
\def \bC{\mathbb C}
\def \bOB{\mathbb {OB}}
\def \bM{\mathbb M}
\def \bOM{\mathbb {OM}}
\def \mA{\mathcal A}
\def \mB{\mathcal B}
\def \mC{\mathcal C}
\def \mR{\mathcal R}
\def \mH{\mathcal H}
\def \mM{\mathcal M}
\def \mV{\mathcal V}
\def \mTOP{\mathcal {TOP}}
\def \mAB{\mathcal {AB}}
\def \bI{\mathbb I}
\def \bK{\mathbb K}
\def \bG{\mathbf G}
\def \bL{\mathbb L}
\def\bN{\mathbb N}
\def\bR{\mathbb R}
\def\bP{\mathbb P}
\def\bZ{\mathbb Z}
\def\bC{\mathbb  C}
\def \bQ{\mathbb Q}
\def\op{\operatorname}

\newcommand{\maru}[1]{\ooalign{
\hfil\resizebox{.8\width}{\height}{#1}\hfil
\crcr
\raise.1ex\hbox{\large$\bigcirc$}}}

\newcommand{\Maru}[1]{\ooalign{
\hfil\resizebox{.6\width}{\height}{#1}\hfil
\crcr
\raise.1ex\hbox{\LARGE$\bigcirc$}}}

\newcommand{\MMaru}[1]{\ooalign{
\hfil\resizebox{.5\width}{\height}{#1}\hfil
\crcr
\raise.1ex\hbox{\Huge$\bigcirc$}}}

\newcommand{\MMMaru}[1]{\ooalign{
\hfil\resizebox{.4\width}{\height}{#1}\hfil
\crcr
\raise.1ex\hbox{\Huge$\bigcirc$}}}

\newcommand{\MMMMaru}[1]{\ooalign{
\hfil\resizebox{.3\width}{\height}{#1}\hfil
\crcr
\raise.1ex\hbox{\Huge$\bigcirc$}}}

\newcommand{\MMMMMaru}[1]{\ooalign{
\hfil\resizebox{.2\width}{\height}{#1}\hfil
\crcr
\raise.1ex\hbox{\Huge$\bigcirc$}}}
%

%%%%%%%%%%%%
\title{Co-operational bivariant theory%\\associated to a contravariant functor
}

\author{Shoji Yokura}

\date{}
\address{Graduate School of Science and Engineering, Kagoshima University, 1-21-35 Korimoto, Kagoshima, 890-0065, Japan}
\email{yokura@sci.kagoshima-u.ac.jp}

\begin{abstract}
For a covariant functor %(with values in the category of graded abelian groups) 
W. Fulton and R. MacPherson defined \emph{an operational bivariant theory} associated to this covariant functor. In this paper we will show that given a contravariant functor %$F^*(-)$ 
one can similarly construct a ``dual" version of an operational bivariant theory, which we call a \emph{co-operational} bivariant theory. %$\bB^{co-op}F^*(X \xrightarrow f Y)$. 
If a given contravariant functor is the usual cohomology theory, then our co-operational bivariant group %$\bB^{co-op}F^*(X \xrightarrow {id_X} X)$ 
for the identity map consists of what are usually called ``cohomology operations". In this sense, our co-operational bivariant theory consists of \emph{``generalized"} cohomology operations.
\end{abstract}

\maketitle


%%%%%%%%%%%%%
\section{Introduction}
In \cite{FM} W. Fulton and R. MacPherson have introduced \emph{bivariant theory} $\mathbb B^*(X \xrightarrow f Y)$  with an aim to deal with Riemann--Roch type theorems for singular spaces and to unify them. $\mathbb B_*(X):= \mathbb B^{-*}(X \xrightarrow {} pt)$ becomes a covariant functor and $\mathbb B^*(X):= \mathbb B^{*}(X \xrightarrow {\op{id}_X} X)$ %becomes 
 a contravariant functor. In this sense  $\mathbb B(X \xrightarrow f Y)$ is called a \emph{bivariant} theory. 
In \S \ref{BT} below we make a quick recall of Fulton--MacPherson's bivariant theory. Here we just remark that in \cite[\S 2.2 Axioms for a bivariant theory]{FM} the value $\mathbb B^*(X \xrightarrow f Y)$ is in the category of graded abelian groups, but as remarked in \cite[Remark, p.22]{FM} $\mathbb B^*(X \xrightarrow f Y)$ can be valued \emph{in an arbitrary category}, e.g., the category of abelian groups (see \cite[\S 6.1.2 Definition of $\mathbb F$, p.60]{FM}), the category of sets (see \cite[\S 4.3 Differentiable Riemann--Roch]{FM}), the category of derived categories (see \cite[\S 7 Grothendieck Duality and Derived Functors]{FM}), etc.
 
As a typical example of such a bivariant theory Fulton and MacPherson defined a bivariant homology theory $\mathbb H^*(X \xrightarrow f Y)$ constructed from the usual cohomology theory $H^*$, which is a multiplicative cohomology theory. Here the multiplicativity is crucial. Given a homology theory, i.e., a covariant functor with values in the category of (graded) abelian groups, in \cite[\S 8.1]{FM} they defined what is called \emph{an operational bivariant theory} or simply \emph{an operational theory}. 
Here we note that in fact there are at least \emph{three more} kinds of operational bivariant theories requiring some other extra conditions, depending on covariant functors to be treated:
\begin{itemize}
\item \cite [\S 8.3]{FM} requiring \emph{one more extra condition}, 
\item \cite[\S 17.1]{Ful} and \cite{AP} requiring \emph{two more extra conditions} and 
\item \cite{GK} requiring \emph{three more extra conditions}.
\end{itemize}
In this paper we do not consider these refined ones for the sake of simplicity.

In this paper, motivated by the definition of an operational bivariant theory given in \cite[\S 8.1]{FM}, we define \emph{a co-operational bivariant theory}\footnote{A co-operational bivariant theory, motivated by the other refined or sophisticated operational bivariant theories listed above will be given and discussed in a different paper.} from \emph{a contravariant functor} which is not necessarily a multiplicative cohomology theory. Let $F^*(-)$ be a contravariant functor such as the usual cohomology theory $H^*(-)$, the K-theory $K(X)$ of complex vector bundles. Then the co-operational bivariant theory $\bB^{coop}F^*(X \xrightarrow f Y)$ is defined in a way analogous to the definition of Fulton--MacPherson's operational bivariant theory associated to a covariant functor $F_*(-)$. In this paper Fulton--MacPherson's operational bivariant theory shall be denoted by $\bB^{op}F^*(X \xrightarrow f Y)$. For example, let $F^*(-)$ be the usual cohomology groups $H^*(-)$. It turns out that for the identity $id_X:X \to X$ the co-operational bivariant group $\bB^{coop}H^*(X \xrightarrow {id_X} X)$ consists of the (ring) homomorphisms $c_g: H^*(X') \to H^*(X')$ for each map $g:X' \to X$ such that for maps 
\begin{equation}\label{maps}
X'' \xrightarrow h X' \xrightarrow g X
\end{equation}
the following diagram commutes:
\begin{equation}\label{cd1}
\xymatrix
{H^*(X'') \ar[d]_{c_{g \circ h}} && H^*(X')\ar[ll]_{h^*} \ar[d]^{c_g}\\
H^*(X'') && H^*(X') \ar[ll]^{h^*}.
}
\end{equation}
In other words, $c_g: H^*(X') \to H^*(X')$ is nothing but what is usually called \emph{a cohomology operation}\footnote{Speaking of ``operation", in \cite{Vishik2} A. Vishik considers an operation between two oriented cohomology theories \cite{LM}(cf. \cite{LP}): An operation from $A^*$ to $B^*$ is a natural transformation from $A^*$ to $B^*$ considered as \emph{contravariant functors} on the category of smooth schemes, i.e., operations commute with \emph{pull-backs (but not necessarily with push-forwards).} On the other hand, in \cite{Sechin} P. Sechin considers \emph{a multiplicative operation} between two oriented cohomology theories.} (e.g., see \cite{Steenrod1, Steenrod2, Steenrod3} and \cite{nLab2}) for the given cohomology group $H^*(-)$. Here we note that we consider such operations for all objects $X'$ over $X$, i.e., $X' \to X$ (in other words, we consider such operations on what is usually called \emph{the category over $X$}), as we define above. In this paper, we can consider a general contravariant functor $F^*(-)$ with values in a certain category. As Fulton and MacPherson considered a homology theory, i.e., a covariant functor with values in the category of graded abelian groups in \cite[\S 8.1]{FM}, in \S \ref{coop} below we consider a contravariant functor $F^*(-)$ with values in the category of graded ablelian groups as a model case. So, in this case, $c_g: H^*(X') \to H^*(X')$ is a homomorphism of graded abelian groups. If $F^*(-)$ is a contravariant functor $F^*(-)$ with values in an arbitrary category, such as the category of sets, abelian groups, derived categories, then  $c_g: H^*(X') \to H^*(X')$ is \emph{a morphism} in the category. Even if we consider such an arbitrary category, a morphism $c_g: H^*(X') \to H^*(X')$  shall be still called \emph{``cohomology" operation}, although it should be called ``a contravariant functor operation" if we ``follow" the way of naming ``cohomology operation". 

A natural transformation $\tau: F^*(-) \to G^*(-)$ between two contravariant functors can be extended to a Grothendieck transformation
$$\tau^{coop}: \bB^{coop}_{\tau}F^*(X \xrightarrow f Y) \to \bB^{coop}G^*(X \xrightarrow f Y)$$
where $\bB^{coop}_{\tau}F^*(X \xrightarrow f Y)$ is a subgroup of $\bB^{coop}F^*(X \xrightarrow f Y)$
and \emph{depends on $\tau$}. 
In particular, for the identity map $id_X:X \to X$, 
$$\bB^{coop}_{\tau}F^*(X \xrightarrow {id_X} X)$$
consists of homomorphisms $c_g:F^*(X') \to F^*(X')$, for all $g:X' \to X$ such that the above diagram (\ref{cd1}) commutes with $H^*(-)$ being replaced by $F^*(-)$ and the following diagram also commutes for some $d \in \bB^{coop}G^*(X \xrightarrow {id_X} X)$:
\begin{equation}\label{cd-FG}
\xymatrix
{ F^*(X') \ar[d]_{c_g } \ar[rr]^{\tau} && G^*(X') \ar[d]^{d_g}  \\
F^*(X') \ar[rr]_{\tau} && G^*(X') \\
}
\end{equation}
And the Grothendieck transformation
$$\tau^{coop}: \bB^{coop}_{\tau}F^*(X \xrightarrow {id_X} X) \to \bB^{coop}G^*(X \xrightarrow {id_X} X)$$
consists of the above commutative diagrams (\ref{cd-FG}). Here we note that obviously the original homomorphism $\tau: F^*(X) \to G^*(X)$ can be considered as the following trivial commutative diagram 
\begin{equation*}%\label{cd-FG-id}
\xymatrix
{ F^*(X) \ar[d]_{id_{F^*(X)}} \ar[rr]^{\tau} && G^*(X) \ar[d]^{id_{G^*(X)}}  \\
F^*(X) \ar[rr]_{\tau} && G^*(X). \\
}
\end{equation*}
Here we give two examples.

Firstly, consider the Chern character $ch: K(-) \to H^{ev}(-;\bQ)$ (e.g., see \cite[\S 3 Properties]{nLab}). Then for each positive integer $k$ we have the Adams operation
$$\Psi^k: K(-) \to K(-)$$
and the Adams-like operation on the even-degree rational cohomology
$$\Psi^k_H:H^{ev}(-;\bQ) \to H^{ev}(-;\bQ)$$
which is defined by $\Psi^k_H(a):=k^r \cdot a$ for each $a \in H^{2r}(-;\bQ)$
and we have the following commutative diagram
\begin{equation}\label{cd3}
\xymatrix
{K(X) \ar[d]_{\Psi^k} \ar[rr]^{ch} && H^{ev}(X;\bQ) \ar[d]^{\Psi^k_H}\\
K(X) \ar[rr]_{ch} && H^{ev}(X;\bQ)
}
\end{equation}
which is natural with respect to the base change, namely we do have the following commutative cube:
\begin{equation}\label{ccube2}
\xymatrix
{K(X') \ar[dd]_{\Psi^k} \ar[rd]^{h^*} \ar[rr]^{ch} && H^{ev}(X';\bQ) \ar'[d][dd]^{\Psi^k_H} \ar[rd]^{h^*} \\
& K(X'') \ar[dd]_{ \Psi^k} \ar[rr]^{ \hspace{1cm} ch}  &&  H^{ev}(X'';\bQ) \ar[dd]^{\Psi^k_H} \\
K(X') \ar'[r] [rr]_{ch \quad \quad }  \ar[rd]_{h^*} && H^{ev}(X';\bQ) \ar[rd]_{h^*} \\
& K(X'') \ar[rr] _{ch \quad \quad  \quad }  &&  H^{ev}(X'';\bQ).
}
\end{equation}
The above commutative diagram (\ref{cd3}) is nothing but \emph{a part} of the homomorphism
$$ch^{coop}: \bB^{coop}_{ch}K^*(X \xrightarrow {id_X} X) \to \bB^{coop}H^*(X \xrightarrow {id_X} X).$$
Here we emphasize that ``a part" means that it is possible that there are many other cohomology operations on $K$-theory and rational even-dimensional cohomology and natural transformations of them like (\ref{cd3}). 
Therefore 
$$ch^{coop}: \bB^{coop}_{ch}K^*(X \to Y) \to \bB^{coop}H^*(X \to Y).$$
is an extension of the above natural transformation (\ref{cd3}) of cohomology operations of $K$-theory and rational even-dimensional cohomology to a Grothendieck transformation.

Secondly, we consider the projection $\overline{pr}: \Omega^*(-) \to CH^*(-)/2$ of the projection from Levine--Morel's algebraic cobordism $\Omega^*(X)$ to the Chow group $CH^*(X)/2$, which is the Chow cohomology group $CH^*(X)/2$ modded out by 2-torsions (see \cite{Vishik}). Then, as shown in \cite{Vishik} (see Brosnan \cite{Bros}, M. Levine \cite{Levine} and A. Merkurjev \cite{Merk})  there exists (unique) operations $S^i:CH^*(X)/2 \to CH^{*+i}(X)/2$, called ``Steenrod operation", such that the following diagram commutes:
\begin{equation}\label{cd-omega}
\xymatrix
{\Omega^*(X) \ar[d]_{S^i_{LN}} \ar[rr]^{\overline{pr}} && CH^*(X)/2 \ar[d]^{S^i}\\
\Omega^*(X) \ar[rr]_{\overline{pr}} && CH^*(X)/2
}
\end{equation}
Here $S^i_{LN}$ is Landwever--Novikov operation (see \cite{LM}) and both Steenrod operation and Landwever--Novikov operations commute with pullbacks, namely we have a commutative cube as the above (\ref{ccube2}), which is nothing but the homomorphism
$$pr^{coop}: \bB^{coop}_{\overline{pr}}\Omega^*(X \xrightarrow {id_X} X) \to \bB^{coop}CH^*(X \xrightarrow {id_X} X)/2.$$
Hence we have the Grothendieck transformation
$$pr^{coop}: \bB^{coop}_{\overline{pr}}\Omega^*(X \to Y) \to \bB^{coop}CH^*(X \to Y)/2$$
is an extension of the above natural transformation (\ref{cd-omega}) of Levine-Morel's algebraic cobordism  and the Chow group $CH^*(X)/2$ to a Grothendieck transformation.

The up-shot is that 
\begin{enumerate}
\item \emph{a co-operational bivariant theory ``captures" or ``interprets" cohomology operations ``bivariant theoretically"}, or \emph{a co-operational bivariant theory is a reasonable general setting to study cohomology operations} and 
\item a Grothendieck transformation
$$\tau^{coop}: \bB^{coop}_{\tau}F^*(X \xrightarrow f Y) \to \bB^{coop}G^*(X \xrightarrow f Y)$$
is an extension or a generalization of a natural transformation of two kind cohomology operations, as shown by the above two examples.
\end{enumerate}

Finally we remark that Fulton--MacPherson's operational bivariant theory is actually described by using the notion of ``homology operation\footnote{In \cite[3.1 Definition]{HM} R. Hardt and C. McCrory define ``stable homology operation" as a natural transformation of homology functor compatible with the suspension isomorphism, just like a stable cohomology operation, which is a natural transformation of cohomology functor compatible with the suspension isomorphism. Hence it is quite natural that ``homology operation" is considered as as a natural transformation of homology functor, just like cohomology operation. In fact, in \cite[Example 4.1.25]{LM} the above Landwever--Novikov operation is described as a homology operation.}" instead of ``cohomology operation". In this sense both theories are really \emph{``operational"} bivariant theories. Furthermore works will be done in a different paper.


%%%%%%%%%%%
\section{Fulton--MacPherson's bivariant theories}\label{BT}

We make a quick review of Fulton--MacPherson's bivariant theory \cite {FM}, since we refer to some axioms required on the theory in later sections.

Let $\mathscr C$ be a category which has a final object $pt$ and on which the fiber product or fiber square is well-defined. Also we consider the following classes:
\begin{enumerate}
\item a class $\mathcal C$ of maps, called ``confined maps" (e.g., proper maps, in algebraic geometry), which are \emph{closed under composition and base change, and contain all the identity maps}, and 
\item a class $\mathcal Ind$ of commutative diagrams, called ``independent squares" (e.g., fiber square, ``Tor-independent" square, in algebraic geometry), 
satisfying that

(i) if the two inside squares in  
$$\CD
X''@> {h'} >> X' @> {g'} >> X \\
@VV {f''}V @VV {f'}V @VV {f}V\\
Y''@>> {h} > Y' @>> {g} > Y \endCD
\quad \quad \qquad \text{or} \qquad \quad \quad 
\CD
X' @>> {h''} > X \\
@V {f'}VV @VV {f}V\\
Y' @>> {h'} > Y \\
@V {g'}VV @VV {g}V \\
Z'  @>> {h} > Z \endCD
$$
are independent, then the outside square is also independent,

(ii) any square of the following forms are independent:
$$
\xymatrix{X \ar[d]_{f} \ar[r]^{\op {id}_X}&  X \ar[d]^f & & X \ar[d]_{\op {id}_X} \ar[r]^f & Y \ar[d]^{\op {id}_Y} \\
Y \ar[r]_{\op {id}_X}  & Y && X \ar[r]_f & Y}
$$
where $f:X \to Y$ is \emph{any} morphism. 
\end{enumerate}

%%%%%%%%%%%
\begin{rem}Given an independent square, 
its transpose is \emph{not necessarily} independent. For example, let us consider the category of topological spaces and continuous maps. Let \emph{any} map be confined, and we allow a fiber square
\quad $\CD
X' @> {g'} >> X \\
@V {f'}VV @VV {f}V\\
Y' @>> {g} > Y \endCD
$ \quad 
to be \emph{independent only if $g$ is proper} (hence $g'$ is also proper). Then its transpose is \emph{not independent unless 
$f$ is proper}. (Note that the pullback of a proper map by any continuous map is proper, because ``proper" is equivalent to ``universally closed", i.e., the pullback by any map is closed.)
\end{rem}
%%%%%%%%%%%%%
%%%%%%%%%%%%%%%%%%%%%%%%
\begin{def}\label{BivariantTheory}
A \emph{bivariant theory} $\mathbb B$ on a category $\mathscr C$ with values in the category of graded abelian groups\footnote{As we mentioned in Introduction, instead of graded abelian groups, we consider also sets, e.g., such as the set of complex structures and the set of Spin structures (see \cite[\S 4.3.2]{FM}),  and categories, e.g., such as the derived (triangulated) category of $f$-perfect complexes (see \cite[\S 7.1 Grothendieck duality]{FM}) as well.}
 is an assignment to each morphism
$ X  \xrightarrow{f} Y$
in the category $\mathscr C$ a graded\footnote{The grading is sometimes ignored. In this case we can consider that the grading is only $0$, i.e, $\bB(X \to Y) = \bB^0(X \to Y)$. For an example for such a case, see \cite[\S 6.1.2 \underline{Definition of} of $\mathbb F$]{FM} where $\mathbb F(X \to Y) = \mathbb F^0(X \to Y)$.}
abelian group $$\bB(X  \xrightarrow{f} Y)$$
which is equipped with the following three basic operations. The $i$-th component of $\bB(X  \xrightarrow{f} Y)$, $i \in \mathbb Z$, is denoted by $\bB^i(X  \xrightarrow{f} Y)$.
\begin{enumerate}
\item {\bf Product}: For morphisms $f: X \to Y$ and $g: Y
\to Z$, the product 
$$\bullet: \bB^i( X  \xrightarrow{f}  Y) \otimes \bB^j( Y  \xrightarrow{g}  Z) \to
\bB^{i+j}( X  \xrightarrow{gf}  Z).$$
%%%%%%%%%%%
\item {\bf Pushforward}: For morphisms $f: X \to Y$
and $g: Y \to Z$ with $f$ \emph {confined}, 

the pushforward 
$f_*: \bB^i( X  \xrightarrow{gf} Z) \to \bB^i( Y  \xrightarrow{g}  Z).$
%%%
\item {\bf Pullback} : For an \emph{independent} square \qquad $\CD
X' @> g' >> X \\
@V f' VV @VV f V\\
Y' @>> g > Y, \endCD
$

the pullback 
$g^* : \bB^i( X  \xrightarrow{f} Y) \to \bB^i( X'  \xrightarrow{f'} Y'). $
\end{enumerate}
%%%%%%%%%%%%%%%%%%%%%%%%

An element $\alp \in \bB(X \xrightarrow f Y)$ is sometimes expressed as follows:
\[
\xymatrix
{
X \ar[rr]_f^{\maru{$\alp$}} && Y
} 
\]
%%%%%%%%%%%%%%%%

These three 
operations are required to satisfy the following seven compatibility 
axioms (\cite [Part I, \S 2.2]{FM}):

\begin{enumerate}
\item[($A_1$)] {\bf Product is associative}: for 
\xymatrix
{
X \ar[r]_f^{\maru{$\alp$}} & Y \ar[r]_g^{\maru{$\be$}} & Z \ar[r]_h^{\maru{$\ga$}} & Z
}
$$(\alp \bullet\be) \bullet \ga = \alp \bullet (\be \bullet \ga).$$
\item[($A_2$)] {\bf Pushforward is functorial}:for 
\xymatrix
{
X  \ar@/^10pt/[rrr]^{\maru{$\alp$}}  \ar[r]_f & Y \ar[r]_g  & Z \ar[r]_h  & W
}
with confined $f, g$, 
$$(g\circ f)_* \alp = g_*(f_*\alp).$$
\item[($A_3$)] {\bf Pullback is functorial}: given independent squares
$$\CD
X''@> {h'} >> X' @> {g'} >> X \\
@VV {f''}V @VV {f'}V @VV {f}V\\
Y''@>> {h} > Y' @>> {g} > Y \endCD
$$
$$(g \circ h)^* = h^* \circ g^*.$$
\item[($A_{12}$)] {\bf Product and pushforward commute}: $f_*(\alp \bullet\be)  = f_*\alp \bullet \be$ \\

for 
\xymatrix
{
X \ar@/^10pt/[rr]^{\maru{$\alp$}}  \ar[r]_f & Y \ar[r]_g  & Z \ar[r]_h^{\maru{$\be$}} & W
}
with confined $f$, 
%%%%%%
\item[($A_{13}$)] {\bf Product and pullback commute}: \quad $h^*(\alp \bullet\be)  = {h'}^*\alp \bullet h^*\be$ \quad for independent squares
 \[
\xymatrix{X' \ar[rrr]^{h''} \ar[d]_{f'}
 &&& X \ar[d]_f^{\maru{$\alp$}}   \\
Y' \ar[d]_{g'}\ar[rrr]^{h'} &&& Y \ar[d]_g^{\maru{$\be$}} \\
Z' \ar[rrr]_h &&& Z
} 
\]
%%%%%%%%%%%
\item[($A_{23}$)] %\label{push-pull} 
{\bf Pushforward and pullback commute}: \quad $f'_*(h^*\alp)  = h^*(f_*\alp)$ \quad for independent squares with $f$ confined 
\[
\xymatrix{
X' \ar[rrr]^{h''} \ar[d]_{f'}  &&& X \ar[d]_f \ar@/^11pt/[dd]^{\maru{$\alp$}} \\
Y' \ar[d]_{g'} \ar[rrr]^{h'} &&& Y \ar[d]_g\\
Z' \ar[rrr]_h &&& Z
} 
\]
%%%%%%%%
\item[($A_{123}$)] {\bf Projection formula}: \, $g'_*(g^*\alp \bullet \be)  = \alp \bullet g_*\be$ \quad for an independent square with $g$ confined 
\[
\xymatrix{
X' \ar[r]^{g'} \ar[d]^{f'}_{\Maru{$g^*\alp$}} 
& X \ar[d]_f^{\maru{$\alp$}}  \\
Y' \ar@/_18pt/[rrrrr]^{\maru{$\beta$}} \ar[r]^g & Y \ar[rrrr]_h^{\quad \Maru{$g_*\be$}} &&&& Z
} 
\]
\end{enumerate}
%%%%%%%%
We also require the theory $\bB$ to have multiplicative units:
\begin{enumerate}
\item[({\bf Units})] For all $X \in \mathscr C$, there is an element $1_X \in \bB^0( X  \xrightarrow{\op {id}_X} X)$ such that $\alp \bullet 1_X = \alp$ for all morphisms $W \to X$ and all $\alp \in \bB(W \to X)$, and such that $1_X \bullet \beta = \beta $ for all morphisms $X \to Y$ and all $\beta \in \bB(X \to Y)$, and such that $g^*1_X = 1_{X'}$ for all $g: X' \to X$.
\end{enumerate}
\end{def}


A bivariant theory unifies both a covariant theory and a contravariant theory in the following sense:
%\begin{def} 
For a bivariant theory $\bB$, its associated covariant functors and contravariant functors are defined as follows:
\begin{enumerate}
\item $\bB_*(X):= \bB(X \xrightarrow {a_X} pt)$ is covariant for confined morphisms and the grading is given by $\bB_i(X):= \bB^{-i}(X \xrightarrow {a_X} pt)$.
\item $\bB^*(X) := \bB(X  \xrightarrow{\op{id}_X}  X)$ is contravariant for all morphisms and the grading is given by $\bB^j(X):= \bB^j(X   \xrightarrow{\op{id}_X}  X)$.
\end{enumerate}
%\end{def}

A typical example of a bivariant theory is the bivariant homology theory $\bH(X \xrightarrow f Y)$ 
constructed from the singular cohomology theory $H^*(-)$, which unifies the singular homology $H_*(X):=\bH^{-*}(X \to pt)$ and the singular cohomology $H^*(X):=\bH^*(X \xrightarrow {\op{id}_X} X)$. Here the underlying category $\mathscr C$ is the category of spaces embeddable as closed subspaces of some Euclidean spaces $\mathbb R^n$ and continuous maps between them (see \cite[\S 3 Topological Theories]{FM}).
More generally, Fulton--MacPherson's (general) bivariant homology theory 
$$h^*(X \to Y)$$
(here, using their notation) 
is constructed from \emph{a multiplicative cohomology theory $h^*(-)$} \cite[\S 3.1]{FM} %(see Definition \ref{biv-hom} below). 
Here the cohomology theory $h^*$ is either ordinary or generalized. %(sometimes, called extra-ordinary). 
A cohomology theory $h^*$ is called \emph{multiplicative} if for pairs $(X,A), (Y,B)$ there is a graded pairing (exterior product)%\footnote{The cup product $\cup: H^i(X, A) \times H^j(X, B) \to H^{i+j}(X, A \cup B)$ is the composite$H^i(X, A) \times H^j(X, B) \xrightarrow {\times} H^{i+j}(X \times X, X \times B \sqcup A \times X) \xrightarrow {\Delta^*} H^{i+j}(X, A \cup B)$ where $\Delta:X \to X \times X$ is the diagomal map.}  
$$h^i(X,A) \times h^j(Y,B) \xrightarrow {\times} h^{i+j}(X \times Y, X \times B \sqcup A \times Y)$$
such that it is associative and graded commutative, i.e., $\alp \times \be = (-1)^{i+j} \be \times \alp$.
A typical example of a multiplicative ordinary cohomology theory is the singular cohomology theory. The topological complex $K$-theory $K(-)$ and cobordism theory $\Omega^*(-)$ are  multiplicative generalized  cohomology theories. We consider the category of spaces embeddable as closed subspaces in some Euclidian spaces $\mathbb R^N$ and continuous maps. For example, \emph{Whitney's embedding theorem} says that any manifold of real dimension $m$ can be embedded as a closed subspace of $\mathbb R^{2m}$. We also note that a complex algebraic variety is embeddable\footnote{This is because the variety $X$ is covered by finitely many affine varieties, which are embedded (as closed subsets) into $\mathbb R^n$ for some $n$, thus it follows from \cite[\S 8.8 Proposition]{Dold} that the variety X is itself embedded (as a closed subset) into $\mathbb R^N$ for some $N$. }  as a closed subspace of some Euclidean space $\mathbb R^N$. We let a confined map be a proper map and an independent square be a fiber square.
%%%%%%%%%
\begin{def}\label{biv-hom}
For a continuous map $f:X \to Y$, choose a map $\phi: X \to \mathbb R^n$ such that $\Phi=(f, \phi): X \to Y \times \mathbb R^n$ defined by $\Phi(x):=(f(x), \phi(x))$ is  a closed embedding\footnote{ Such a map always exist, since our space is embeddable as a closed subspace of some $\mathbb R^N$, thus this embedding is considered as $\phi:X \to \mathbb R^N$, then $\Phi=(f, \phi)$ is also a closed embedding.}. Then we define
\begin{equation}\label{embed}
h^*(X \to Y):= h^{i+n}(Y \times \mathbb R^n, Y \times \mathbb R^n \setminus \Phi(X)).
\end{equation}
\end{def}
%%%%%%%%%
\begin{thm}(\cite[p.34-p.38]{FM}) The above definition (\ref{embed}) is independent of the choice of the embedding $\phi$, thus $\Phi$, and $h^*(X \to Y)$ is a bivariant theory.
\end{thm}
%%%%%%%%%%%%%
\begin{rem} 
\begin{enumerate}
\item By the definition (\ref{embed}) we have
$h^i(X \xrightarrow {\op{id}_X} X) = h^i(X).$
Indeed, since $\op{id}_X: X \to X$ is obviously a closed map (embedding), we can choose $\phi:X \to pt$ so that
$\Phi=(\op{id}_X, \phi):X \to X \times pt \cong X$ is a closed map. Hence %by the definition (\ref{embed}) 
we have 
$$h^i(X \xrightarrow {\op{id}_X} X)= h^i(X, (X \times pt) \setminus \Phi(X))=h^i(X, \emptyset)=h^i(X).$$
\item $h^{-i}(X \xrightarrow {a_X} pt) = h^{n-i}(\mathbb R^n, \mathbb R^n \setminus \Phi(X))=:h^{n-i}(\mathbb R^n, \mathbb R^n \setminus X)$ where $\Phi=(a_X, \phi):X \to pt \times \mathbb R^n =\mathbb R^n$ is a closed embedding. If $h^*=H^*$ is the singular cohomology, then $h^{-i}(X \xrightarrow {a_X} pt) =h^{n-i}(\mathbb R^n, \mathbb R^n \setminus X) =H^{n-i}(\mathbb R^n, \mathbb R^n \setminus X) =:H^{BM}_i(X)$ is the Borel--More homology group (e.g., see \cite{Ful}, \cite[B.1]{PS}).
\end{enumerate}
\end{rem}
%%%%%%%%%%%%%%
\begin{defn}
(\cite[\S 2.7 Grothendieck transformation]{FM}) \label{groth}
Let $\bB, \bB'$ be two bivariant theories on a category $\mathscr C$. 
A {\it Grothendieck transformation} from $\bB$ to $\bB'$, $\ga : \bB \to \bB'$
is a collection of homomorphisms
$\bB(X \to Y) \to \bB'(X \to Y)$
for a morphism $X \to Y$ in the category $\mathscr C$, which preserves the above three basic operations: 
\begin{enumerate}
\item $\ga (\alp \bullet_{\bB} \be) = \ga (\alp) \bullet _{\bB'} \ga (\be)$, 
\item $\ga(f_{*}\alp) = f_*\ga (\alp)$, and 
\item $\ga (g^* \alp) = g^* \ga (\alp)$. 
\end{enumerate}
\end{defn}
%%%%%%%%%%%%%%%%%%%%%%%%%%%%%

A Grothendieck transformation $\ga: \bB \to \bB'$ induces natural transformations $\ga_*: \bB_* \to \bB_*'$ and $\ga^*: \bB^* \to {\bB'}^*$.
 
%%%%%%%%%
\begin{rem}\label{contra}
Let $t: h^* \to \widetilde h^*$ be a natural transformation of two multiplicative cohomology theory. Then we get the associated Grothendieck transformation
\begin{equation}\label{bivari-h}
t:h^*( X \xrightarrow f Y) \to \widetilde h^*( X \xrightarrow f Y) 
\end{equation}
since we have $t: h^{*+n}(Y \times \mathbb R^n, Y \times \mathbb R^n \setminus \Phi(X)) \to \widetilde h^*(Y \times \mathbb R^n, Y \times \mathbb R^n \setminus \Phi(X))$. For example, the Chern character
$ch: K^0(-) \to H^*(-) \otimes \mathbb Q$ induces the Grothendieck transformation
$ch: K^0(X \xrightarrow f Y) \to H^*(X \xrightarrow f Y) \otimes \mathbb Q$ (see \cite[\S 3.2.2 \underline{Examples}]{FM}).
\end{rem}
%%%%%%%%%%%
\section{Operational bivariant theory}\label{obt}
Given a covariant functor or a homology theory, Fulton and MacPherson have defined what is called \emph{an operational bivariant theory} or \emph{an operational theory} \cite[\S 8 \underline{Operational Theories}]{FM} (also, see \cite[\S 17. 1 and \S 17.2]{Ful}). As we remarked in Introduction, in this paper we consider the operational bivariant theory defined in \cite[\S 8.1]{FM}. The case of the other refined operational bivariant theories will be considered in a different paper where we define
a co-operational bivariant theory ``corresponding" to these refined operational bivariant theories.

Let $h_*$ be a homology theory, i.e., a covariant functor with values in graded abelian groups such that the functorial (pushforward) homomorphism $f_*:h_*(X) \to h_*(Y)$ is defined for a confined map $f:X \to Y$.  Then its \emph{operational bivariant theory}, denoted by $\bB^{op}h_*(X \xrightarrow f Y)$, is defined as follows. For a map $f:X \to Y$, an element $c \in \bB^{op}h_*^i(X \xrightarrow f Y)$ is defined to be a collection of
homomorphisms\footnote{In \cite[\S 8 \underline{Operational Theories}]{FM}, $c_g$ is denoted by $c_{Y'}$ and in \cite[\S 17. 1]{Ful} $c_g$ is denoted by $c^{(k)}_g$.}
\begin{equation*}
c_g:h_k(Y') \to h_{k-i}(X')
\end{equation*}
for all $k \in \mathbb Z$, all $g:Y' \to Y$ and the fiber square \, \, \, 
$$\CD
X' @> {g'}  >> X\\
@V {f'} VV @VV fV\\
Y' @>> {g} > Y.\\
\endCD
$$
These homomorphisms $c_g$ are
required to be \emph{compatible with pushforward} (for confined maps), i.e., for a fiber diagram
\, \, \, 
$$\CD
X'' @> {k'}  >> X' @> {g'}  >> X\\
@V {f''} VV @VV f'V @VV fV\\
Y'' @>> {k} > Y' @>> {g} > Y.\\
\endCD
$$
where $k$ is confined and thus $k'$ is confined as well, 
the following diagram commutes
\begin{equation}
\xymatrix
{h_*(X'')  \ar[rr]^{k'_*} && h_*(X')\\
h_*(Y'') \ar[u]^{c_{g\circ k}}  \ar[rr]_{k_*} && h_*(Y') \ar[u]_{c_g}
}
\end{equation}

If the category $\mathscr C$ has a final object $pt$ and $h_*(pt)$ has a distinguished element $1$, then the homomorphism $ev:\bB^{op}h_*(X \to pt) \to h_*(X)$ defined by $ev(c) := (c(\op{id}_{pt}))(1)$ is called the evaluation homomorphism.
Let $\bB$ be a bivariant theory. Then its \emph{associated operational bivariant theory} $\bB^{op}$ is defined to be the operational bivariant theory constructed from the covariant functor $\bB_*(X) = \bB(X \to pt)$. Then we have the following canonical Grothendieck transformation
$$ op:\bB \to \bB^{op}$$
defined by, for each $\alp \in \bB(X \to Y)$,
$$
op(\alp):= \{ (g^*\alp) \bullet : \bB_*(Y') \to \bB_*(X') | g:Y' \to Y \}.
$$

In this case it is not clear whether one could construct a Grothendieck transformation of the associated operational bivariant theories from a natural transformation of two covariant functors. To be more precise, if $t :h_*(-) \to \widetilde h_*(-)$ is a natural transformation of two covariant functors, then it is not clear whether one could construct a Grothendieck transformation $t:\bB^{op}h_*(X \xrightarrow f Y) \to \bB^{op}\widetilde h_*(X \xrightarrow f Y)$, which is an ``operational bivariant theoretic analogue" of (\ref{bivari-h}). A kind of similar problem is discussed in \cite[\S 8.2]{FM}. Suppose that $\bB$ is a bivariant theory, $h_*$ is a covariant functor and there are homomorphisms $\phi(X): \bB_*(X)=\bB(X \to pt) \to h_*(X)$, covariant for confined maps, and taking $1 \in \bB^*(pt)=\bB_*(pt)$ to $1 \in h_*(pt)$. Then a question is whether there exists a unique Grothendieck transformation $\Phi:\bB(X \to Y) \to \bB^{op}h_*(X \to Y)$ such that the associated map $\Phi(X): \bB_*(X) \to \bB^{op}h_*(X \to Y)$ followed by the evaluation map $ev_X:\bB^{op}h_*(X \to pt) \to h_*(X)$, i.e., $ev_X \circ \Phi(X)$, is equal to the given homomorphism $\phi(X):\bB_*(X) \to h_*(X)$. The answer to this question is negative, however a positive answer to a modified question is affirmative, details for which see \cite{Yokura-ope} (cf. \cite{BSY-jpaa}). 
%%%%%%%%
\begin{rem} For example, if we consider the refined operational bivariant theory associated to the Chow homology theory \cite[Proposition 17.3.1]{Ful}, then we have $\bB^{op}h^{-p}(X \to pt) \cong h_p(X)$. We will discuss more details in a different paper where we define
co-operational bivariant theories ``corresponding" to other refined operational bivariant theories.
\end{rem}

%%%%%%%%%%%%%
\section{Co-operational bivariant theory}\label{coop}
In this section, motivated by the above Fulton--MacPherson's operational bivariant theory, we consider another kind of operational bivariant theory associated to \emph{a contravariant functor} which may not be a multiplicative cohomology theory. This new bivariant theory shall be called \emph{a co-operational bivariant theory}.

If a contravariant functor $h^*$ is a multiplicative cohomology theory and $t: h^* \to \widetilde h^*$ is a natural transformation between two multiplicative cohomology theories, then as recalled above, one can extend $t$ to a Grothendieck transformation $t:h^*(X \to Y) \to \widetilde h^*(X \to Y)$ between their associated bivariant theories in such a way that for an identity map $id_X:X \to X$, $t:h^*(X \xrightarrow {id_X} X) \to \widetilde h^*(X \xrightarrow {id_X} X))$ is equal to the original natural transformation $t:h^*(X) \to \widetilde h^*(X)$. If it is not multiplicative, one cannot make such a construction. For example, the Chern class $c: K^0(-) \to H^*(-;\bZ)$ is not a natural transformation between two multiplicative cohomology theories, as proved by B. Totaro \cite{Totaro} (cf. \cite{Groth} and \cite{Segal}).
Our co-operational bivariant theory is defined as follows:
\begin{thm}[``Co-operational" bivariant theory]
Let $F^*(-)$ be a contravariant functor with values in the category of graded abelian groups, defined on a category $C$ with all fiber squares independent.
For each %confined\footnote{Since we use Gysin pushforward homomorphisms $c_g: F^k(X') \to F^{k+i}(Y')$ for a map $f':X' \to Y'$ in the definition of co-operational bivariant group, we require $f':X' \to Y'$ to be a confined map. If $f:X \to Y$ is confined, then the pullback $f':X' \to Y'$ is also confined.} 
map $f:X \to Y$ we set
$$ c \in \mathbb B^{coop}F^{* i}(X \xrightarrow f Y)$$
to be a collection of %Gysin \emph{pushforward} 
homomorphisms
$$c_g: F^k(X') \to F^{k+i}(Y')$$
for all $k \in \mathbb Z$, all $g:Y' \to Y$  and the fiber square \, \, \, 
$$\CD
X' @> {g'}  >> X\\
@V {f'} VV @VV fV\\
Y' @>> {g} > Y.\\
\endCD
$$ \, \, \, 

And these homomorphisms $c_g$ are
required to be compatible with pullback, i.e., for a fiber diagram
\begin{equation}\label{cd-0}
\CD
X'' @> {h'}  >> X' @> {g'}  >> X\\
@V {f''} VV @VV f'V @VV fV\\
Y'' @>> {h} > Y' @>> {g} > Y.\\
\endCD
\end{equation}
the following diagram commutes
\begin{equation}\label{cd-01}
\xymatrix
{F^k(X'') \ar[d]_{c_{g \circ h}} && F^k(X')\ar[ll]_{(h')^*} \ar[d]^{c_g}\\
F^{k+i}(Y'') && F^{k+i}(Y') \ar[ll]^{h^*}.
}
\end{equation}
We define the following bivariant-theoretic operations:
\begin{enumerate}
%%%%%%%%%%
\item {\bf Product}: The product
$$\bullet: \mathbb B^{coop}F^{*i}(X \xrightarrow f Y) \otimes \mathbb B^{coop}F^{*j}(Y \xrightarrow g Z) \to \mathbb B^{coop}F^{*i+j}(X \xrightarrow {g \circ f} Z)$$
is defined by, for $c \in \mathbb B^{coop}F^{*i}(X \xrightarrow f Y) $ and $d \in \mathbb B^{coop}F^{*j}(Y \xrightarrow g Z) $,
$$(c \bullet d)_h:= d_h \circ c_{h'}: F^k(X') \xrightarrow {c_{h'}}  F^{k+i}(Y') \xrightarrow {d_h} F^{k+i+j}(Z')$$
Here we consider the following fiber squares:
\begin{equation}\label{cd-1}
\CD
X' @> {h''} >> X \\
@V {f'}VV @VV {f}V\\
Y' @> {h'} >> Y \\
@V {g'}VV @VV {g}V \\
Z'  @>> {h} > Z \endCD 
\end{equation}
%%%%%%
\item {\bf Pushforward}: For $X \xrightarrow f Y \xrightarrow g Z$ %with $f$ being confined (hence $g \circ f$ being confined as well)
$$f_*: \mathbb B^{coop}F^{*i}(X \xrightarrow {g \circ f} Z) \to \mathbb B^{coop}F^{* i}(Y \xrightarrow g Z)$$
is defined by, for $c \in \mathbb B^{coop}F^{*i}(X \xrightarrow {g \circ f} Z)$ 
$$(f_*c)_h := c_h \circ (f')^*: F^k(Y') \xrightarrow{(f')^*} F^k(X')  \xrightarrow{c_h} F^{k+i}(Z').$$
Here we use the above commutative diagram (\ref{cd-1}).
%%%%%%%
\item {\bf Pullback}: For a fiber square 
$$\CD
X' @> {g'}  >> X\\
@V {f'} VV @VV fV\\
Y' @>> {g} > Y\\
\endCD
$$
%where $f$ is confined,
$$g^*: \mathbb B^{coop}F^{* i}(X \xrightarrow f Y ) \to \mathbb B^{coop}F^{* i}(X' \xrightarrow {f'} Y' )$$
is defined by, for $c \in \mathbb B^{coop}F^{*i}(X \xrightarrow f Y) $
$$(g^*c)_h := c_{g \circ h}: F^k(X'') \to F^k(Y'').$$
Here we use the above commutative diagrams (\ref{cd-0}).
\end{enumerate} 
Then $\mathbb B^{coop}F^{* i}(X \xrightarrow f Y)$ %(defined on all confined maps $f: X \to Y$) 
becomes a bivariant theory, i.e., satisfies seven axioms of bivariant theory, listed in \S \ref{BT}.
\end{thm}
%%%%%%%%%%
\begin{proof} The proofs %are `basically' straightforward, thus they 
could be omitted. However, since the definition of our co-operational bivariant theory is different from that of Fulton--MacPherson's operational bivariant theory, for the sake of convenience or time-saving for the reader, we write down proofs, giving just comments for trivial ones. %Doing this would probably give the reader some idea for a proof for that Fulton--MacPherson's operational bivariant theory is a bivariant theory as well.
\begin{enumerate}
\item First we remark that these three operations are compatible with pullbacks. Here we show only the compatibility of the pushforward. For which we consider the following fiber squares:
\begin{equation}\label{cd-1-3}
\CD
X'' @> {k''} >>X' @> {h''} >> X \\
@V {f''}VV @V {f'}VV @VV {f}V\\
Y'' @> {k'} >> Y' @> {h'} >> Y \\
@V {g''}VV @V {g'}VV @VV {g}V \\
Z'' @>> {k} >Z'  @>> {h} > Z \endCD 
\end{equation}
and consider the following diagrams:
\begin{equation}
\xymatrix
{F^k(X'') \ar[dd]_{c_{h\circ k}} &&&& F^k(X') \ar[llll]_{(k'')^*} \ar[dd]^{c_h}\\
& F^k(Y'') \ar[lu]_{(f'')^*} \ar[dl]^{(f'')_*c_{h \circ k}}  && F^k(Y')\ar[ll]_{(k')^*} \ar[dr]_{(f')_*c_h} \ar[ur]^{(f')^*} &\\
F^{k+i}(Z'') &&&& F^{k+i}(Z') \ar[llll]^{k^*}
}
\end{equation}
The outer square, the upper trapezoid, the left and right triangles are commutative, therefore the lower trapezoid is commutative, which means that the pushforward  is compatible with pullback.

\item {\bf $(A_1)$ Product is associative}: It is because the product is defined as composition of homomorphisms and composition is associative.
\item  {\bf $(A_2)$ Pushforward is functorial}:for 
\xymatrix
{
X  \ar@/^10pt/[rrr]^{\maru{$c$}}  \ar[r]_f & Y \ar[r]_g  & Z \ar[r]_h  & W
}


$$(g\circ f)_* c= (g_* \circ f_*)c, \quad c \in \mathbb B^{coop}F^{* i}(X \xrightarrow {h \circ g \circ f} Y).$$
Consider the following fiber squares:
\begin{equation}\label{cd-2}
\CD
X' @> {j'''} >> X \\
@V {f'}VV @VV {f}V\\
Y' @> {j''} >> Y \\
@V {g'}VV @VV {g}V \\
Z'  @> {j'} >> Z \\
@V {g'}VV @VV {g}V \\
W'  @>> {j} > W 
\endCD 
\end{equation}
Then by the definition we have
\begin{align*}
((g\circ f)_*c)_j & := c_j \circ (g'\circ f')^* , \quad c_j \in F^k(X') \to F^{k+i}(W') \\
                  & = c_j \circ ((f')^* \circ (g')^*) \\
                  & = (c_j \circ (f')^*) \circ (g')^*\\ 
                  & = (f_*c)_j \circ (g')^*\\  
                  & = (g_*(f_*c))_j \\
                  & = ((g_* \circ f_*)c)_j: F^k(Z') \to  F^{k+i}(W')          
\end{align*}
Therefore $(g\circ f)_* c= (g_* \circ f_*)c \in \mathbb B^{coop}F^{* i}(Z \xrightarrow g W).$

\item {\bf $(A_3)$ Pullback is functorial}: It is clear by the definition of pullback.
\item {\bf $(A_{12})$ Product and pushforward commute}: Let us consider the above diagrams (\ref{cd-2}): Let $\alp \in 
\mathbb B^{coop}F^{* i}(X \xrightarrow {g \circ f} Y)$ and $\beta \in \mathbb B^{coop}F^{* j}(Z \xrightarrow h W).$
\begin{align*}
\left ( f_*(\alp \bullet \beta)\right)_j & = (\alp \bullet \beta)_j \circ (f')^*\\
& = (\beta_j \circ \alp_{j'} ) \circ (f')^*\\
& = \beta_j \circ (\alp_{j'} \circ (f')^*)\\
& = \beta_j \circ (f_*\alp)_{j'}\\
& = ((f_*\alp) \bullet \beta)_j
\end{align*}
Hence we have $f_*(\alp \bullet \beta) = (f_*\alp) \bullet \beta.$
\item {\bf $(A_{13})$ Product and pullback commute}: Let us consider the above diagrams (\ref{cd-1-3}): Let $\alp \in 
\mathbb B^{coop}F^{* i}(X \xrightarrow f Y)$ and $\beta \in \mathbb B^{coop}F^{* j}(Y \xrightarrow g Z).$
\begin{align*}
\left ( h^*(\alp \bullet \beta)\right)_k & = (\alp \bullet \beta)_{h \circ k}\\
& = \beta_{h \circ k} \circ \alp_{h' \circ k'}\\
& = (h^*\beta)_k \circ ((h')^*\alp)_{k'}\\
& = ((h')^*\alp) \bullet h^*\beta)_k
\end{align*}
Hence we have $h^*(\alp \bullet \beta) = (h')^*\alp \bullet h^*\beta.$
\item {\bf $(A_{23})$ Pushforward and pullback commute}: Let us consider the above diagrams (\ref{cd-1-3}): Let $\alp \in 
\mathbb B^{coop}F^{* i}(X \xrightarrow {g \circ f} Z)$.
\begin{align*}
\left ((f')_*(h^*\alp) \right)_k & = (h^*\alp)_k \circ (f'')^*\\
& = \alp_{h \circ k} \circ (f'')^*\\
& = (f_*\alp)_{h \circ k} \\
& = (h^*(f_*\alp))_k
\end{align*}
Hence we have $(f')_*(h^*\alp)= h^*(f_*\alp).$
\item {\bf $(A_{123})$ Projection formula}: 
Let us consider the following diagrams:
\begin{equation}
\xymatrix
{
& X'\ar[rr]^{\qquad \qquad g'} \ar'[d][dd]_{f'} && X \ar[dd]^f &&\\
\widetilde {X'} \ar[ur]^{\widetilde {k''}} \ar[dd]_{\widetilde{f'}}  \ar[rr]^{\qquad \qquad \widetilde {g'}} && \widetilde X \ar[dd] \ar[ur]_{\widetilde {k'}} & & &\\
& Y' \ar'[r][rr]^g && Y \ar[rr]^h && Z\\
\widetilde {Y'} \ar[ur]^{k''} \ar[rr]_{\widetilde g} && \widetilde Y \ar[ur]_{k'}^{\widetilde f \quad } \ar[rr]_{\widetilde h} && \tilde Z \ar[ur]_k
}
\end{equation}
Let $\alp \in  \mathbb B^{coop}F^{* i}(X \xrightarrow f Y)$ and $\beta \in  \mathbb B^{coop}F^{* j}(Y' \xrightarrow {h \circ g} Z)$
\begin{align*}
\left ((g')_*(g^*\alp \bullet \beta ) \right)_k & = (g^*\alp \bullet \beta )_k \circ (\widetilde {g'})^*\\
& = (\beta_k \circ (g^*\alp)_{k''} ) \circ (\widetilde {g'})^*\\
& = \beta_k \circ  ((g^*\alp)_{k''} \circ (\widetilde {g'})^* )\\
& = \beta_k \circ  ((\widetilde g)^* \circ \alp_{k'} )\\
& = ( \beta_k \circ  (\widetilde g)^*) \circ \alp_{k'} \\
& = (g_*\beta)_k \circ \alp_{k'} \\
& = (\alp \bullet g_*\beta)_k
\end{align*}
Hence we have $(g')_*(g^*\alp \bullet \beta )= \alp \bullet g_*\beta.$

\end{enumerate}
\end{proof}
%%%%%%%%
\begin{rem} In Fulton--MacPherson's bivariant theory $\bB^*$ the bivariant pushforward $f_*:\bB^*(X \xrightarrow {g \circ f} Z) \to \bB^*(Y \xrightarrow g Z)$ is defined for a confined map $f:X \to Y$. As one can see, in our co-operational bivariant theory, for the bivariant pushforward $f_*: \mathbb B^{coop}F^{*i}(X \xrightarrow {g \circ f} Z) \to \mathbb B^{coop}F^{* i}(Y \xrightarrow g Z)$ we do not need the confined-ness of the map $f:X \to Y$.
\end{rem}
%%%%%%%
\begin{rem} Let $\mathbb B^*(X \to Y)$ be a bivariant theory and let $c \in \mathbb B^i(X \xrightarrow f Y)$.
Then for any morphism $g:Y' \to Y$ and consider the fiber square
$$\CD
X' @> {g'}  >> X\\
@V {f'} VV @VV fV\\
Y' @>> {g} > Y.\\
\endCD
$$
Then $(g^*c)^*: B_k(Y'):=\mathbb B^{-k}(Y' \to pt) \to  \mathbb B^{-k+i}(X' \to pt) =:B_{k-i}(X')$
is defined by $(g^*c)^*(a):=g^*c \bullet a$ (as defined in \cite[\S 2.5 Gysin homomorphisms]{FM}). For fiber squares
$$\CD
X'' @> {h'}  >> X' @> {g'}  >> X\\
@V {f''} VV @VV f'V @VV fV\\
Y'' @>> {h} > Y' @>> {g} > Y.\\
\endCD
$$ \, \, \, 
the following diagram commutes:
$$
\xymatrix
{B_k(X'') \ar[rr]^{(h')_*} && B_k(X') \\
 B_{k-i}(Y'') \ar[u]^{((hg)^*c)^*}  \ar[rr]_{h_*} && B_{k-i}(Y') \ar[u]_{(g^*c)^*}.
}
$$ 
Then considering the collection $\{ g^*c | g:Y' \to Y\}$ gives us an operational bivariant theory from the given bivariant theory. So, conversely Fulton--MacPherson's operational bivariant theory is one obtained from a covariant functor. Similarly, our co-operational bivariant theory is motivated by considering the following:
$(g^*c)_*: B^k(X'):=\mathbb B^k(X' \xrightarrow {id} X') \to \mathbb B^{k+i}(Y' \xrightarrow {id} Y') =: B^{k+i}(Y')$
is defined by $(g^*c)_*(b):=f'_*(b \bullet g^*c)$ (as defined in \cite[\S 2.5 Gysin homomorphisms]{FM}). Here we remark that, as pointed out in the first paragraph of \cite[\S 2.5 Gysin homomorphisms, p.26]{FM}, when we consider the Gysin homomorphism $\theta_*$ induced by a bivariant element $\theta \in \bB (X \xrightarrow f Y)$ the map $f$ is assumed to be confined. For fiber squares
$$\CD
X'' @> {h'}  >> X' @> {g'}  >> X\\
@V {f''} VV @VV f'V @VV fV\\
Y'' @>> {h} > Y' @>> {g} > Y.\\
\endCD
$$ \, \, \, 
the following diagram commutes:
$$
\xymatrix
{B^k(X'') \ar[d]_{((g\circ h)^*c)_*} && B^k(X')\ar[ll]_{(h')^*} \ar[d]^{(g^*c)_*}\\
 B^{k+i}(Y'') && B^{k+i}(Y') \ar[ll]^{h^*}.
}
$$  
\end{rem}
%%%%%%%%%%%%%%
\begin{pro} Let $\mathbb B$ be a bivariant theory. Let $B^*(X):=\mathbb B^*(X \xrightarrow {id_X} X)$ be the associated contravariant functor. Then, restricted to confined morphisms $f:X \to Y$, there exists a canonical Grothendieck transformation
$$\ga^{coop}: \mathbb B^i(X \to Y) \to \mathbb B^{coop}B^{* i}(X \xrightarrow f Y )$$ 
which is defined by
$$\ga^{coop} (\alp) := \left \{ (g^*\alp)_*:B^*(X') \to B^*(Y') \, | \, \, g:Y' \to Y \right \}$$
Here $(g^*\alp)_*$ is defined by $(g^*\alp)_*(b)= f'_*(b \bullet g^*\alp)$, where $b \in B^*(X')=\mathbb B^*(X' \xrightarrow {id_{X'}} X')$ and $g^*\alp \in \mathbb B^*(X' \xrightarrow {f'} Y')$ and we consider the following square
$$\CD
X' @> {g'}  >> X\\
@V {f'} VV @VV fV\\
Y' @>> {g} > Y.\\
\endCD
$$
Then we have
\begin{enumerate}
\item $\ga^{coop} (\alp \bullet \be)= \ga^{coop}(\alp) \bullet \ga^{coop}(\be)$.
\item $\ga^{coop} (f_*\alp)= f_*\ga^{coop}(\alp)$.
\item $\ga^{coop} (f^*\alp)= f^*\ga^{coop}(\alp)$.
\end{enumerate}
\end{pro}
%%%%%%%%

Suppose that we have natural transformation of two contravariant functors:
$$T:F^* \to G^*.$$
There does not necessarily exist a Grothendieck transformation
$$\ga_T: \mathbb B^{coop} F^*(X \to Y)  \to  \mathbb B^{coop} G^*(X \to Y)  .$$
For that we need to modify the source bivariant theory $\mathbb B^{coop} F^*(X \to Y)$ slightly as follows.
\begin{defn} We define 
$$\mathbb B^{coop}_T F^*(X \to Y) \quad (\subset \mathbb B^{coop}F^{* i}(X \xrightarrow f Y) )$$
to be a collection of %Gysin \emph{pushforward} 
homomorphisms
$$c_g: F^k(X') \to G^{k+i}(Y')$$
for all $k \in \mathbb Z$, all $g:Y' \to Y$  and the fiber square \, \, \, 
\begin{equation}\label{cd-6}
\CD
X' @> {g'}  >> X\\
@V {f'} VV @VV fV\\
Y' @>> {g} > Y.\\
\endCD
\end{equation}
such that
\begin{enumerate}
\item these homomorphisms $c_g$ are
required to be compatible with pullback, i.e., for a fiber diagram
\begin{equation*}
\CD
X'' @> {h'}  >> X' @> {g'}  >> X\\
@V {f''} VV @VV f'V @VV fV\\
Y'' @>> {h} > Y' @>> {g} > Y.\\
\endCD
\end{equation*}
the following diagram commutes
\begin{equation*}
\xymatrix
{F^k(X'') \ar[d]_{c_{g \circ h}} && F^k(X')\ar[ll]_{(h')^*} \ar[d]^{c_g}\\
F^{k+i}(Y'') && F^{k+i}(Y') \ar[ll]^{h^*}.
}
\end{equation*}

\item for all $k \in \mathbb Z$, all $g:Y' \to Y$  and the above fiber square (\ref{cd-6}) there exists a %Gysin pushforward 
homomorphism $d_g:  G^k(X') \to G^{k+i}(Y')$ such that 
\begin{enumerate}
\item $d \in  \mathbb B^{coop} G^*(X \to Y)$ and 
\item the following diagram commutes:
\begin{equation}\label{cd-7}
\xymatrix
{F^k(X') \ar[d]_{c_g} \ar[rr]^T && G^k(X') \ar[d]^{d_g}\\
F^{k+i}(Y') \ar[rr]_T && G^{k+i}(Y'). 
}
\end{equation}
\end{enumerate}
\end{enumerate}
\end{defn}
%%%%%%%%%%
\begin{rem}
We emphasize that therefore we automatically have the following commutative cube:
\begin{equation}\label{ccube-2}
\xymatrix
{ F^*(X') \ar[dd]_{c_g } \ar[rd]^{h^*} \ar[rr]^{T} && G^*(X') \ar'[d][dd]^{d_g} \ar[rd]^{h^*} \\
& F^*(X'') \ar[dd]_{ c_{g \circ h}} \ar[rr]^{ \hspace{1cm} T}  &&  G^*(X'') \ar[dd]^{d_{g \circ h}} \\
F^*(X') \ar'[r] [rr]_{T \quad \quad }  \ar[rd]_{h^*} && G^*(X') \ar[rd]_{h^*} \\
& F^*(X'') \ar[rr] _{T \quad \quad  \quad }  &&  G^*(X'').
}
\end{equation}
\end{rem}

\begin{cor}
Then a Grothendieck transformation which we are looking for is 
$$\ga_T: \mathbb B^{coop}_T F^*(X \to Y)  \to  \mathbb B^{coop} G^*(X \to Y)$$
which is a collection of the above commutative diagram (\ref{cd-7}).
\end{cor}

%%%%%%%%%
%%%%%%%%%
\section{$ \mathbb B^{coop}F^{*}(X \xrightarrow {id_X} X)$ and cohomology operations}

$ \mathbb B^{coop}F^{*}(X \xrightarrow {id_X} X)$ consists of a collection of homomorphisms
$$c_g: F^k(X') \to F^{k+i}(X')$$
for all $g:X' \to X$  and for $X'' \xrightarrow {h} X' \xrightarrow g X$, the following diagram commutes:
$$
\xymatrix
{F^k(X'') \ar[d]_{c_{g \circ h}^k} && F^k(X')\ar[ll]_{(h')^*} \ar[d]^{c_g^k}\\
F^{k+i}(X'') && F^{k+i}(X') \ar[ll]^{h^*}.
}
$$ 
In other words, $c_g: F^k(X') \to F^{k+i}(X')$ is what is called \emph{a cohomology operation}, which is \emph{a natural self-transformation} of a contravariant functor. Therefore $ \mathbb B^{coop}F^{*}(X \xrightarrow {id_X} X)$ consists of natural self-transformations of the contravariant functor $F^*$. The above commutative cube (\ref{ccube-2}) means that $T$ is \emph{a natural transformation between two contravariant functors equipped with natural self - transformations}.
%%%%%%%
\begin{rem} Corresponding to the identity $id_X:X \to X$, $c_g: F^k(X') \to F^{k+i}(X')$ is a cohomology operation \emph{which is associated to the identity $id_X:X \to X$}. So, for a general map $f: X \to Y$,
$c_g: F^k(X') \to G^{k+i}(Y')$ could be called \emph{a ``generalized" cohomology operation associated to the map $f:X \to Y$}, where we consider
the fiber square \, \, \, 
\begin{equation*}
\CD
X' @> {g'}  >> X\\
@V {f'} VV @VV fV\\
Y' @>> {g} > Y.\\
\endCD
\end{equation*}
\end{rem}
%%%%%
\begin{rem}
As we remarked above, if $F^*(-)$ has values in an arbitrary category, then $ \mathbb B^{coop}F^{*}(X \xrightarrow {id_X} X)$ should be a collection of \emph{morphisms}, hence, for example, if we consider the category of rings, then it is a collection of ring homomorphisms; if we consider the category of sets, then it is a collection of maps. Even if we consider the category of rings or (graded) abelian groups, if we ignore some structures, e.g., ignore all the structures and consider only just sets, then it has to be a collection of maps. In other words, it depends on what kind of situations we consider. 
\end{rem}
%%%%%%%
%%%%%%%%
\begin{rem} As pointed out by T. Annala (in private communication), if we consider a more refined co-operational bivariant theory $ \widehat{\mathbb B}^{coop}F^{*}$ requiring certain extra conditions, then we will obtain an isomorphism $\widehat{\mathbb B}^{coop}F^{*}(X \xrightarrow {id_X} X) \cong F^*(X)$, which is similar to the case of the refined operational bivariant theory of Chow homology theory treated in \cite[\S 17.3]{Ful}.
Details of this and some other refined co-operational bivariant theories will be treated in a different paper.
\end{rem}
%%%%%%%%%

Here are some examples.
\begin{ex}\label{Adam} For this example, e.g., see Karoubi's book \cite[IV, 7.13 Theorem and V, 3.37 Theorem]{Kar}.
The $k$-th Adams operation 
$$\Psi^k:K(X) \to K(X)$$
 for the topological K-theory is a ring homomorphism and it turns out that the Adams operation is the only operation satisfying a ring homomorphism\footnote{If we consider only the additive structure, i.e., a group structure, ignoring the product structure, then there may be other operations.}. The Adams-like operation 
$$\Psi^k_H:H^{ev}(X;\bQ) \to H^{ev}(X;\bQ)$$
 for the even cohomology with rational coefficients is defined by
 $\Psi^k_H(x) := k^r x$ for $x \in H^{2r}(X, \bQ)$. The Chern character $ch:K(X) \to H^{ev}(X ;\bQ)$ is compatible with these Adams operations, i.e., the following diagram commutes:
$$
\xymatrix
{
K(X) \ar[d]_{\Psi^k}\ar[rr]^{ch}  && H^{ev}(X;\bQ) \ar[d]^{\Psi^k_H}\\
K(X) \ar[rr]_{ch} && H^{ev}(X;\bQ)
}
$$
which is compatible with pullback, i.e., for any continuous map $h:X \to Y$ we have the commutative cube:
\begin{equation*}
\xymatrix
{ K(Y) \ar[dd]_{\Psi^k } \ar[rd]^{h^*} \ar[rr]^{ch} && H^{ev}(Y;\bQ) \ar'[d][dd]^{\Psi^k_H} \ar[rd]^{h^*} \\
& K(X) \ar[dd]_{\Psi^k} \ar[rr]^{ \hspace{1cm} ch}  &&  H^{ev}(X;\bQ)\ar[dd]^{\Psi^k_H} \\
K(Y)\ar'[r] [rr]_{ch \quad \quad }  \ar[rd]_{h^*} && H^{ev}(Y;\bQ)\ar[rd]_{h^*} \\
& K(X) \ar[rr] _{ch \quad \quad  \quad }  &&  H^{ev}(X;\bQ).
}
\end{equation*}
 From the Chern character $ch:K(X) \to H^{ev}(X ;\bQ)$ we get the following Grothendieck transformation
$$\ga_{ch}: \mathbb B^{coop}_{ch} K^*(X \to Y)  \to  \mathbb B^{coop} H^{ev *}(X \to Y)_{\mathbb Q}.$$
Here we note that in this case $\mathbb B^{coop}_{ch} K^*(X \xrightarrow {id_X} X) = \mathbb B^{coop} K^*(X \xrightarrow {id_X} X)$, which consists of only Adams operations, provided that we require the homomorphism $c_g:K(X') \to K(X')$ to be a ring homomorphism. 
\end{ex}
%%%%%%%%%%%%%%%
\begin{ex} Let $\Omega^*(X)$ be Levine--Morel's algebraic cobordism on smooth schemes and let $pr: \Omega^*(X) \to CH^*(X)$ be the canonical map to the Chow cohomology defined by $pr([V \xrightarrow h X]):=h_*[V] \in CH^*(X)$. Let $LN_I: \Omega^*(X) \to \Omega^{* + |I|}(X)$ be the Landwever--Novikov operation for a partition $I = (n_1 \geq n_2 \geq \cdots \geq n_i>0)$ with $|I|:=n_1 + n_2 + \cdots + n_i$ being the degree of $I$, and $LN^{CH}_I: CH^*(X) \to CH^*(X)$ be the Chow cohomology version of the Landwever--Novikov operation (see \cite[\S 3]{Zain}). Then we have the following commutative diagram:
$$
\xymatrix
{
\Omega^*(X)  \ar[d]_{LN_I}\ar[rr]^{pr}  && CH^*(X) \ar[d]^{LN^{CH}_I}\\
\Omega^*(X)\ar[rr]_{pr} && CH^*(X)
}
$$
From the canonical map $pr: \Omega^*(X) \to CH^*(X)$ we get the following Grothendieck transformation
$$\ga_{pr}: \mathbb B^{coop}_{pr} \Omega^*(X \to Y)  \to  \mathbb B^{coop} CH^*(X \to Y).$$
The Landwever--Novikov operations $LN_I: \Omega^*(X) \to \Omega^*(X)$ become elements of $\mathbb B^{coop}_{pr} \Omega^*(X \xrightarrow {id_X} X)$.  
\end{ex}
\begin{ex}\label{LN} Let $CH^*(X)/2$ be the Chow cohomology mod out by the 2-torsion part and
let $\overline{pr}:\Omega^*(X) \to CH^*(X)/2$ be the composition of the above $pr:\Omega^*(X) \to CH^*(X)$ and the projection $CH^*(X) \to CH^*(X)/2$. Then Brosnan, Levine and Merkurjev (see \cite{Vishik} and \cite{Bros}, \cite{Levine} and \cite{Merk} ) showed that there exists (unique) operation $S^i:CH^*(X)/2 \to CH^*(X)/2$, called Steenrod operation, making the following diagram commutative
$$
\xymatrix
{
\Omega^*(X)  \ar[d]_{S^i_{LN}}\ar[rr]^{\overline{pr}}  && CH^*(X)/2 \ar[d]^{S^i}\\
\Omega^*(X)\ar[rr]_{\overline{pr}} && CH^*(X)/2
}
$$
which is compatible with pullback. Here $S^i_{LN} := LN_I$ where  $I = (n_1 \geq n_2 \geq \cdots \geq n_i>0)$ with $n_1 =n_2 = \cdots =n_i=1$, thus $|I|=i$.\\\\
\end{ex}


%\begin{qu}At the moment I have been unable to find a reasonable map from $\mathbb B^{co-op}H^{*}(X \xrightarrow {id_X} X)$ to $H^*(X)$. What is a reasonable map?
%\end{qu}

\noindent
{\bf Acknoweldgements}:

This work was motivated by a question posed by Toru Ohmoto  (in private communication) after his recent work \cite{Ohmoto} on solving Kazarian's conjecures \cite{Kaz} on Thom polynomials. The author thanks him for posing the question and some useful discussion. The author also thanks Burt Totaro for answering some questions about his work \cite{Totaro} and Toni Annala for some comments on the first draft of this paper.
Some of this work was done during the author's visiting Fakult\"at f\"ur Mathematik, Universit\"at Wien, in May 2023. The author thanks Herwig Hauser for a very nice hospitality and financial support. 
The author is supported by JSPS KAKENHI Grant Numbers JP19K03468 and JP23K03117.\\



\begin{thebibliography}{ams}

\bibitem{AP} D. Anderson and S. Payne, {\it Operational $K$-theory},  Doc. Math., {\bf 20} (2015), 357--399.


\bibitem{Bros} P. Brosnan, {\it Steenrod operations in Chow theory}
, Trans. Amer. Math. Soc. 355 (2003), no.5, 1869--1903.

\bibitem{BSY-jpaa}  J.-P. Brasselet, J. Sch\"urmann and S. Yokura, {\it On Grothendieck transformations in Fulton--MacPherson's bivariant theory}, J. Pure Appl. Algebra, {\bf 211} (2007), 665--684.

\bibitem{Dold} A. Dold, {\it Lectures on Algebraic Topology}, 2nd Edition, Springer--Verlag, Berlin Heidelberg 2407 New York 1980.

\bibitem{Ful}
W. Fulton, {\it Intersection theory}, Springer Verlag (1984).
  
\bibitem{FM}
W. Fulton and R. MacPherson, {\it Categorical framework for the study of  
  singular spaces},  Mem. Amer. math. Soc., {\bf 243} (1981).
  
\bibitem{GK}J. L. Gonz\'alez and K. Karu, {\it Bivariant algebraic cobordism}, Algebra \& Number Theory, {\bf 9} , No.6 (2015), 1293--1336.

\bibitem{Groth} A. Grothendieck, {\it La  th\'eorie des classes de Chern}, Bull. Soc. Math. France, {\bf 86} (1958), 137--154.

\bibitem{HM} R. Hardt and C. McCrory, {\it Steenrod operations in subanalytic homology}, Compos. Math., {\bf 39}, No. 3 (1979), 333--371.

\bibitem{Kar} M. Karoubi, {\it K-Theory; An introduction}, Grundlehren der mathematischen Wissenschaften 226, Springer 1978. 
 
\bibitem{Kaz} M. E. Kazarian, {\it Multisingularities, cobordisms, and enumerative geometry}, Russian Math. Surveys, {\bf 58}, Issue 4 (2003), 665--724.

\bibitem{Levine}
M. Levine, {\it Steenrod operations, degree formulas and algebraic cobordism}, Pure and Applied Mathematical Quaterly Journal, volume in honor of the 60th birthday of Robert MacPherson, {\bf 3}, No. 1 (2007) 283--306.

\bibitem{LM}
M. Levine and F. Morel,
{\it Algebraic Cobordism},
Springer Monographs in Mathematics, Springer-Verlag (2007).

\bibitem{LP}
M. Levine and R. Pandharipande,
{\it Algebraic cobordism revisited},
Invent. Math., {\bf 176} (2009), 63-130.

\bibitem{Merk} A. Merkurjev, {\it Steenrod operations and degree formulas},  J. Reine Angew. Math. {\bf 565} (2003), 13--26.

\bibitem{nLab} nCatLab, {\it Adams operation in nLab}, 
https://ncatlab.org/nlab/show/Adams+operation.

\bibitem{nLab2} nCatLab, {\it Cohomology operation in nLab}, 
https://ncatlab.org/nlab/show/cohomology+operation.

%\bibitem{Ohmoto-1} T. Ohmoto, private communication,  November 24, 2022.

\bibitem{Ohmoto} T. Ohmoto, {\it Multi-singularities and Hilbert's 15th problem}, Lecture Notes (in Japanese) for the conference ``Singularities of Differentiable Maps and Its Applications -- In commemoration of Goo Ishikawa's all years of hard work--", at Hokkaido University, March 6 - 8, 2023.

\bibitem{PS} C. A. M. Peters and J. H. M. Steenbrink, {\it Mixed Hodge Structures}, Springer--Verlag, 2008. 

\bibitem{Sechin} P. Sechin, {\it On the structure of algebraic cobordism}, Adv. Math., {\bf 333}(2018), 314--349.

\bibitem{Segal} G. Segal, {\it The multiplicative group of classical cohomology}, Q.J. Math. Oxf. {\bf 26} (1975), 289--293.

\bibitem{Steenrod1} N. E. Steenrod, {\it Cohomology Operations and Obstructions to Extending Continuous Functions}, Colloquium Lectures, the Annual Summer Meeting of the Amer. Math. Soc. at Pennsylvania State University in August, 1957.

\bibitem{Steenrod2} N. E. Steenrod, {\it Cohomology Operations, and Obstructions to Extending Continuous Functions}, Adv. Math., {\bf 8} (1972), 371--416.

\bibitem{Steenrod3} N. E. Steenrod and D. B. A. Epstein, {\it Cohomology Operations}, Annals of Mathematics Studies, Princeton University Press, 1962.
https://people.math.rochester.edu/faculty/doug/otherpapers/steenrod-epstein.pdf

\bibitem{Totaro} B. Totaro, {\it The total Chern class is not a map of multiplicative cohomology theories}, 
Math. Z., {\bf 212} (1993), 527--532.

\bibitem{Vishik} A. Vishik, {\it Algebraic cobordism, Landwever--Novikov and Steenrod operations, symmetric operations}, in School and Conference on Algebraic $K$-theory and its Applications, The Abdus Salam ICTP, 2007.

\bibitem{Vishik2} A. Vishik, {\it Stable and Unstable operations in Algebraic Cobordism},  Annales Scientifiques de l'Ecole Normale Superieure, {\bf 52}, No.3 (2019), 561--630.

\bibitem{Yokura-ope} S. Yokura, {\it Bivariant theories of constructible functions and Grothendieck transformations}, Topology Appl.,  {\bf 123} (2002), 283--296.

\bibitem{Zain} K. Zainoulline, {\it Localized operations on $T$-equivariant oriented cohomology of projective homogeneous varieties}, J. Pure Appli. Algebra, {\bf 226}, Issue 3 (2022), 106856.


\end{thebibliography}

\end{document}